\documentclass{report}
\usepackage{lmodern}
\usepackage{graphicx} % Required for inserting images
\usepackage[a4paper, left=2cm, right=2cm, top=3cm, bottom=3cm]{geometry}
\usepackage{amsmath} %usa questo per scrivere le formule matematiche
\usepackage{parskip} %Questo pack elimina l'indentazione fastidiosa 
\usepackage{amsthm} % Questo pack permette di scrivere in modo carino i teoremi,
\usepackage{amssymb}
\usepackage{xcolor}    % Per i colori
\usepackage{framed}    % Per le cornici
\usepackage{mdframed}  % Cornici più avanzate, definizioni, proposizioni ecc.
\usepackage{xfrac} %per fare le diagonali carine

%Quelli che seguono sono i risultati dell'uso del pack amsthm
\newtheorem{theo}{Theorem}
\newtheorem{prop}{Proposition}
\newtheorem{crlr}{Corollario}[theo] % Numerato in base ai teoremi
\newtheorem{defi}{Definition}
\newtheorem{note}{Note}
\newtheorem{xpl}{Example}
\newenvironment{pop}{}{}
\renewenvironment{proof}{{\noindent\bfseries Proof. }}{\qed}



\title{Algebra}
\author{Fabio Colletti}
\date{September 2026}

\begin{document}

\maketitle

\chapter{Groups}
\section{Introduction}
    \begin{defi}[Group]
        A \textbf{Group} $(G,\cdot)$ is a set $G$ which has an operation $\cdot : G \times G \rightarrow G $ that is:
        \begin{enumerate}
            \item $f\cdot (g \cdot h) = (f \cdot g) \cdot h \quad \forall f,g,h \in G$ \quad \textbf{(Associativity)}
            \item $\exists \space e \in G \quad \text{t.c.} \quad g \cdot e = e \cdot g = g \quad \forall g \in G$ \quad \textbf{(Neutral Element $\mathbf{e}$)}
            \item $\forall \space g \in G \space \exists g^{-1} \quad \text{t.c.} \quad g^{-1}\cdot g = g \cdot g^{-1} = e $ \quad \textbf{(Inverse Element)}
        \end{enumerate}
    \end{defi}\vspace{.5 cm}
    \begin{note}
        There's the Associativity property, so I can always write:
        $$f\cdot g\cdot h=(f\cdot g)\cdot h = f\cdot(g \cdot h)$$
        I'll also obmit the dot when it is clear in the context.
    \end{note}\vspace{.5 cm}
    \begin{defi}
        A group is called:
        \begin{enumerate}
            \item An \textbf{Abelian Group} if $f\cdot g = g\cdot f \quad \forall f,g \in G$
            \item A \textbf{Finite Group} if $|G| < \infty$
        \end{enumerate}
    \end{defi}\vspace{.5 cm}
    \begin{defi}[Subgroup]
        Let $(G,\cdot)$ be a group. Let $H \subseteq G$ be a subset.\\
        $H$ is a $G$\textbf{-Subgroup} if $(H,\cdot)$ is a group.
    \end{defi}\vspace{.5 cm}
    \begin{prop}
        Let $G$ be a group.\\ 
        Let $H\quad t.c.\quad  |H|<\infty $ a G subset. It is a $G$-subgroup if it is closed in the group operation. 
    \end{prop}
    \begin{proof}
        We need to demonstrate that $(H,\cdot)$ is a group.\\
        Let's start by assuming that $H$ is finite and closed in the group operation.
        \begin{itemize}
            \item $\forall g \in H, \quad g \in G \quad \Rightarrow \text{Associativity is granted.}$
            \item $|H| < \infty\quad\Rightarrow grad(g)=n_g \quad\Rightarrow g^{n_g}=e \quad\Rightarrow e\in H,\quad g^{-1}=g^{n_g-1}\in H$
        \end{itemize}
        
    \end{proof}\vspace{.5cm}
    \newpage    
    
\section{Lagrange Theorem}
    \begin{defi}[Congruent in module]
        Let $G$ be a group, $H\leq G$ a subgroup, $f,g\in G$.\\
        $f$ is called \textbf{congruent to $g$ mod $H$} if: 
        $$f^{-1}\cdot g \quad\in H  $$
        \\In this case I'll write: $f\equiv g \quad (mod H)\quad$ or $\quad f \equiv_{H} g$
    \end{defi}\vspace{.5 cm}
    \begin{prop}
        The property of being congruent in module is a class of equivalence.\\
        \begin{proof}
            Exercise.
        \end{proof}
    \end{prop}\vspace{.5 cm}
    \begin{prop}[Class of equivalence]
        Let $(G,\cdot)$ be a group, $H\le G$ a subgroup, $g \in G$.\\ 
        Therefore the class of equivalence of $g$ in $H$ is:
        $$gH:=\{g\cdot h \quad|\quad h\in H  \}$$
        \begin{proof}
            Exercise.
        \end{proof}
    \end{prop}\vspace{.5 cm}
    \begin{defi}
        The set of $gH$ for all $g \in G$ is written: $\sfrac{G}{H}$
    \end{defi}\vspace{.5 cm}
    \begin{defi}
        let a group be $G$ and $H\leq G$ a subgroup. $H$ is called \textbf{Normal} if: $\quad gH = Hg \quad\forall g \in G$.\\ In that case It's written: $\quad H\unlhd G$. 
    \end{defi}\vspace{.5 cm}
    
\end{document}